% Generated draft: requires mathematical and source review.
\documentclass[11pt]{article}
\usepackage[margin=25mm]{geometry}
\usepackage{fontspec}
\setmainfont{Latin Modern Roman}
\usepackage{amsmath,amssymb,mathtools,unicode-math}
\setmathfont{Latin Modern Math}
\usepackage{enumitem,needspace,xurl,hyperref}
\setlength{\parindent}{0pt}
\setlength{\parskip}{5pt}
\setlength{\emergencystretch}{3em}
\newcommand{\N}{\mathbb{N}}
\newcommand{\Z}{\mathbb{Z}}
\newcommand{\Q}{\mathbb{Q}}
\newcommand{\R}{\mathbb{R}}
\newcommand{\C}{\mathbb{C}}
\newcommand{\F}{\mathbb{F}}
\newcommand{\A}{\mathbb{A}}
\newcommand{\PP}{\mathbb{P}}
\newcommand{\E}{\mathbb{E}}
\newcommand{\eps}{\varepsilon}
\newcommand{\dd}{\,\mathrm{d}}
\begin{document}
\section*{UnsolvedMath ID: 800676}
\textbf{Source problem code:} catalogue entry
\subsection{Definitions and mathematical statement}
Fix an integer \(d\geq 2\). Retain the source notation \(\mathscr{B}_d\) for its parameter space of degree-\(d\) foliations. The supplied record does not define this parameter space, its complex analytic structure, the underlying ambient manifold, or the convention for degree. The formulation below is therefore relative to those missing source definitions.

Write
\[
\mathscr{B}_d^{c}=\{\mathcal{F}\in\mathscr{B}_d: \mathcal{F}\text{ admits a nonzero directed positive closed current}\}.
\]
Here a positive closed current is a positive current \(T\) satisfying \(\mathrm{d}T=0\); directedness means tangency to the foliation on its regular locus, subject to the source's intended current type and behavior at singularities. Those details are not specified in the supplied record. Nonzero is the usual necessary convention: allowing the zero current would make the existence condition vacuous.

The question is whether, for every integer \(d\geq2\), there exist proper complex analytic subsets \(A_{d,n}\subsetneq\mathscr{B}_d\), indexed by \(n\in\N\), such that
\[
\mathscr{B}_d^{c}=\bigcup_{n\in\N}A_{d,n}.
\]
Analyticity and properness are taken in the intended parameter space \(\mathscr{B}_d\).

For context, the source defines \(\mathscr{H}_d\subseteq\mathscr{B}_d\) as the set of foliations with only hyperbolic singularities and describes it as open and real Zariski dense. It states that \(\mathscr{B}_d^{a}\cap\mathscr{H}_d\) is analytic, citing [FS3] and a bound on the degree of the relevant curve when all singularities are hyperbolic. The supplied record does not define \(\mathscr{B}_d^{a}\), the relevant curves, or hyperbolicity. It also states that \(\mathscr{B}_d^{c}\) is closed and attributes the equality
\[
\mathscr{B}_d^{c}\cap\mathscr{H}_d
=
\mathscr{B}_d^{a}\cap\mathscr{H}_d
\]
to part of Theorem 3.1. These are supplied contextual assertions; their hypotheses and supporting statements are not reproduced in the record.
\subsection{Short English statement}
For each degree \(d\geq2\), ask whether the foliations admitting a nonzero directed positive closed current form a countable union of proper analytic subsets of the intended parameter space. The record labels the problem open. Its missing ambient definitions and current conventions require source clarification before this becomes a fully specified mathematical statement.
\subsection{Sources}
\begin{itemize}
\item Supplied local UnsolvedMath registry record: permanent ID \(800676\), source problem code \texttt{catalogue entry}, title \emph{Holomorphic Foliation Theory --- Problem 3}; AMR Open Problem Lists collection. The supplied status is \texttt{open}, not an independently verified literature assessment.
\item Dinh, Sibony, \emph{Some Open Problems on Holomorphic Foliation Theory} (2018), Problem 3. Supplied source URL: \url{https://arxiv.org/abs/1807.05640}. The record attributes the problem to Dinh and Sibony, with proposed year 2018.
\item Supplied provenance reports access on 2026-10-16 and extraction by \texttt{arxiv\_latex}. This access date is supplied metadata; no independent source visit or verification is claimed.
\item The source excerpt cites [FS3] for analyticity of \(\mathscr{B}_d^{a}\cap\mathscr{H}_d\), but supplies no bibliographic entry identifying [FS3]. It also refers to Theorem 3.1 without reproducing its complete statement or hypotheses.
\end{itemize}
\end{document}
